\section*{Part II: The Anatomy of Segments}

\subsection*{13. A Segment Is a Causal Hypothesis About Behavior}

To segment a market is to propose that certain differences between people explain differences in the decisions they make. The categories matter only insofar as they help predict buying, usage, outcome, cost, or retention.

A table of ages and income levels may be useful, but it is not automatically strategic. Consider 2 27-year-old software engineers in the same city, with similar compensation. One has just been promoted to technical lead and must design an AI feature; the other enjoys mathematics but has no near-term project. The first may urgently need review and accountability. The second may prefer a textbook or an open-source project. They look alike demographically but inhabit different commercial situations.

A useful segment should have relatively low variation in the decisive buying and delivery factors. That does not require all members to be identical. It requires similarity in what matters.

\subsection*{14. The 7 Practical Axes of Segmentation}

There are many ways to divide a market. Start with these 7 and ask which are actually causal:

\begin{table*}[t]
\centering
\small
\begin{tabularx}{\textwidth}{l p{5cm} X}
\toprule
\textbf{Axis} & \textbf{What It Distinguishes} & \textbf{Potential Use} \\
\midrule
Situation and trigger & What changed, and why now? & A system release, hiring mandate, compliance deadline, or career transition \\
Job and desired outcome & What progress is sought? & Deploy a feature, certify staff, reduce scrap, or earn an industry credential \\
Capability and readiness & What can the user already do? & Beginner vs.\ working engineer; small vs.\ large client implementation capacity \\
Economics and authority & Who pays and approves? & Self-pay, employer pay, departmental budget, tender, or retailer purchase \\
Constraints & What makes adoption hard? & Schedule, geography, security, regulatory obligations, or language \\
Channel and social context & How is the decision reached? & Peer community, retail shelf, procurement, or campus administration \\
Usage and service intensity & How much support or repetition is needed? & Self-service, guided cohort, or bespoke enterprise implementation \\
\bottomrule
\end{tabularx}
\caption{The 7 practical axes of segmentation.}
\label{tab:ch7-segmentation-axes}
\end{table*}

Frequently the best segment uses 2 or 3 dimensions together: a buyer with a particular job and trigger operating under a specific constraint. Using 7 simultaneously can overfit a tiny audience and make the segment impossible to reach.

\subsection*{15. Segment by What Predicts a Different Product Decision}

A segmentation variable is useful when changing its value changes a meaningful business choice: curriculum, product architecture, price, acquisition channel, sales proof, staffing, service level, or success measurement.

If all segments receive the same offer, purchase through the same channel, require the same proof, and cost the same to serve, the segmentation may be decorative. Conversely, 1 variable may split the business in 2. An individual learner paying tuition personally and a corporate training manager allocating departmental budget often require entirely different offers even when both ultimately teach the same technical topic.

A simple test: \textbf{What would we do differently for this group?} If the answer is nothing, reconsider why it is called a strategic segment.

\subsection*{16. An Underserved Segment Is Not Merely an Unhappy Segment}

Customers complain about almost everything. Some frustrations are cheap irritations rather than reasons to switch. An underserved group is one for whom existing feasible alternatives perform materially below the group's required outcome, given its time, money, and risk constraints.

It helps to separate 3 cases:

\begin{itemize}
    \item \textbf{Underperformance:} Existing solutions repeatedly fail at an outcome this segment cares about enough to pursue.
    \item \textbf{Over-service:} Incumbents provide expensive complexity or service the group does not value; a simpler, cheaper option may be better.
    \item \textbf{Nonconsumption:} People have a meaningful job but cannot practically access or afford existing solutions, so they use workarounds or do nothing.
\end{itemize}

Clayton Christensen's theory of disruptive innovation is particularly relevant to some over-served and non-consuming footholds, but not every niche entrant is a disruptor [R3]. A specialist premium product can also enter from the top, where incumbents fail to provide unusually demanding service. The point is to identify the actual mismatch, not to claim every small venture is changing an industry.

\subsection*{17. How to Find the Underserved in Customer Language}

Ask people about the last time their current solution failed. What did they have to do themselves? Which compromises did they accept? Who else became involved? Which tasks were excluded from the vendor's scope? What did they pay to repair the outcome? Did they search again? Did they stop trying?

Look for phrases attached to observable behavior: ``We had to ask another team to review every response''; ``the program was too basic for our experienced engineers''; ``the college couldn't schedule the lab because no instructor was available''; ``we built a spreadsheet to finish what the system couldn't do.''

A complaint is a clue. A \textbf{costly workaround, attempted purchase, delayed project, or paid substitute} is stronger evidence of a possible opening.

\subsection*{18. Explore the Abandoned Edges of Successful Products}

Dominant products tend to optimize for their largest or most profitable groups. That creates natural edges. Large standardized courses may prioritize content scale over deep review. Enterprise software may optimize for large accounts with administrators. A widely sold FMCG product may prioritize mass taste, pack size, and retail convenience over a small use occasion.

Study whom the incumbent rejects, prices out, fails to onboard, or burdens with undesirable features. Study people who use only 1 small part of a bundled product. Study users who maintain parallel spreadsheets or manual processes despite paying for software. Study buyers whose existing solution is technically sufficient but procedurally impossible to approve.

The attractive opening is often a customer whose job is \textbf{incompatible with the incumbent's chosen operating model}, not a feature the incumbent has never heard of.

\subsection*{19. The Problem Must Be Acute to Someone with Authority}

A team may experience severe pain but lack a budget. A manager may have budget but not prioritize the problem. A consumer may experience pain but face an attractive free substitute. A student may strongly value education but be unable to protect time.

For a beachhead to support early sales, the urgency must connect to the decision-making system. Distinguish end-user urgency, economic-buyer urgency, and organizational permission. For an enterprise provider, winning the end user's affection does not necessarily unlock procurement. For a D2C institute, a learner's anxiety does not necessarily convert into self-financed tuition.

This is why the statement ``our customer desperately needs it'' must always be followed by: \emph{Who can do something about that need now?}

\subsection*{20. Frequency and Severity Produce Different Beachheads}

A problem that occurs every day but costs a few minutes may support low-priced, recurring software. A problem that appears once a year but threatens a multimillion-rupee compliance event may support a costly expert engagement. A course purchase may happen once but demand 150 hours of effort. A snack purchase may happen weekly with 2 seconds of deliberation.

Do not multiply frequency and severity into a magic desirability score. Ask instead about the appropriate economic form: annual contract, one-off project, per-use charge, tuition, subscription, replenishment, or sponsorship. The initial segment must support a transaction with an economics and buying rhythm the founding team can manage.
